\documentclass{amsart}
\usepackage{amsthm}
\usepackage{graphicx,textcomp}
\usepackage{amsmath,bm,amsfonts,mathrsfs,amssymb}
\usepackage{hyperref}
\newtheorem{theorem}{Theorem}
\newtheorem{lemma}[theorem]{Lemma}
\newtheorem{prop}[theorem]{Proposition}
\newtheorem{defi}[theorem]{Definition}
\newtheorem{rem}[theorem]{Remark}
\newtheorem{cor}[theorem]{Corollary}
%\usepackage{comment,refcheck}
\usepackage{appendix}
\usepackage[normalem]{ulem}
\usepackage{graphicx,color}
\usepackage{xcolor}
\def\blue{\color{blue}}
\def\red{\color{red}}

\counterwithout{equation}{section}   

\begin{document}
\title[Determinant of the Dirichlet-to-Neumann map]{The determinant of the Dirichlet-to-Neumann map for a surface with boundary and periods of holomorphic differentials on its double}

\author{Dmitrii Korikov}
	\address{St. Petersburg Department of Steklov Mathematical Institute
of Russian Academy of Sciences, 27 Fontanka, St. Petersburg, Russia, \url{https://orcid.org/0000-0002-3212-5874}}
	\email{dmitrii.v.korikov@gmail.com}
	
	\author{Alexey Kokotov}
	\address{Department of Mathematics \& Statistics, Concordia University, 1455 De Maisonneuve Blvd. W. Montreal, QC  H3G 1M8, Canada, \url{https://orcid.org/0000-0003-1940-0306}}
	\email{alexey.kokotov@concordia.ca}
	
\begin{abstract}
Let $(M,g)$ be a smooth orientable $2d$ Riemannian manifold of genus $\mathfrak{g}$ with  Riemannian metric $g$ and connected boundary $\Gamma$. Let $\Lambda$ be the Dirichlet-to-Neumann map on $\Gamma$ and let ${\rm det}_\zeta(\Lambda)$ be its (modified, i. e. with  zero mode excluded) $\zeta$-regularized determinant. It is well-known that the quantity ${\rm det}_\zeta(\Lambda)/|\Gamma|$ (where $|\Gamma|$ is the length of $\Gamma$) is a conformal invariant. It was shown by Edward and Wu (\cite{EV}) that this invariant equals one for $\mathfrak{g}=0$; in the case $\mathfrak{g}>0$ Guillarmou and Guillop\'e \cite{Guillarmou} found two explicit expressions for this invariant through the Rouelle and (respectively) the Selberg zeta-functions of the two surfaces of negative constant curvature from the conformal class of $(M,g)$: one is  of infinite volume and complete whereas another has geodesic boundary. 

We present an elementary counterpart of the formulae of Guillarmou and Guillop\'e using  the periods of holomorphic differentials on the double $2M$ of $M$ only.  Our approach is based on the properties of the Hilbert transform of $M$ \cite{B,HilbKor} and the Kontsevich-Vishik-Friedlander-Guillemin regularization of the determinants of pseudodifferenial operators \cite{KV,Ww,F}. In particular, a connection between the length spectra of (uniformized) $M$, $2M$  and the periods of holomorphic differentials on $2M$  is established.
\end{abstract}

	\maketitle

\section*{Introduction}
Let $(M,g)$ be a (smooth) orientable surface of genus $\mathfrak{g}$ with metric $g$ and connected boundary $\Gamma$. We assume that $\Gamma$ is parametrized by the arc length and identify it with the boundary of the disk $\mathbb{D}$ of radius $|\Gamma|/2\pi$. The Dirichlet-to-Neumann (DN  in the sequel) map of $(M,g)$ is the self-adjoint operator in $L_2(\Gamma,dl_g;\mathbb{C})$ given by $\Lambda f:=\partial_\nu u^f|_\Gamma$, where $u^f$ is the harmonic extension of $f\in H^{1}(\Gamma;\mathbb{C})\equiv {\rm Dom}(\Lambda)$ into $(M,g)$ and $\nu$ is the unit exterior normal to $\Gamma$. Let $\partial_\gamma$ denotes differentiation with respect to the length along $\Gamma$. It is well-known \cite{LeeU} that $\Lambda$ is a $\Psi$DO of order one, admitting decomposition
\begin{equation}
\label{DN PDO decomp}
\Lambda=\Lambda_0+\tilde{\Lambda}, \qquad \Lambda_0=|\partial_\gamma|,
\end{equation}
where $\tilde{\Lambda}$ is smoothing $\Psi$DO. Note that $\Lambda_0$ is the DN map of $\mathbb{D}$.

In the sequel we restrict all the operators to the Hilbert space $\mathscr{H}:={\rm Ran}(\Lambda)=L_2(\Gamma,dl_g;\mathbb{C})\ominus\mathbb{C}$, thereby making $\Lambda$ and $\partial_\gamma$ invertible. As follows from formula (4.2.62), \cite{Grubb} ("no pure $\log$-term in the heat asymptotics")  the operator $\zeta$-function $\zeta_{\Lambda}(s)$ is regular at $s=0$. The (modified) $\zeta$-regularized determinant of $\Lambda$ is defined as ${\rm det}_\zeta(\Lambda):={\rm exp}\big(-\zeta'_\Lambda(0)\big)$; according to  \cite{Guillarmou} it is related to the length spectrum of $(M\backslash\Gamma,h_\infty)$ (where $h_\infty$ is the unique complete hyperbolic metric on $M\backslash\Gamma$) via
\begin{equation}
\label{GG formula}
{\rm det}_\zeta(\Lambda)/|\Gamma|=\Big((2\pi s)^{-\mathfrak{g}}\mathscr{R}(s)\Big)\Big|_{s=0}/(1-\mathfrak{g})
\end{equation}
for $\mathfrak{g}>1$, where $\mathscr{R}(s):=\prod_{\gamma}(1-e^{-s\,|\gamma|_h})$ is the Ruelle zeta-function, the product is over all primitive closed geodesics in $(2M,h_\infty)$ (note that the left-hand side is a confromal invariant). (There is a similar relation also established in  \cite{Guillarmou} with Selberg zeta-function  of the hyperbolic surface with geodesic boundary conformally equivalent to $M$, we do not quote this result in full for brevity sake.)


For the trivial case $\mathfrak{g}=0$, the explicit formulas $\zeta_{\Lambda_0}(s)=2\zeta(s)\Big(\frac{2\pi}{|\Gamma|}\Big)^{-s}$, $\zeta_{\partial_\gamma}(s)=\zeta_{\Lambda_0}(s){\rm cos}\Big(\frac{\pi s}{2}\Big)$ and $\zeta(0)=-\frac{1}{2}$, $\zeta'(0)=-\frac{{\rm log}(2\pi)}{2}$ (where $\zeta$ is the Riemann zeta function) imply
\begin{equation}
\label{GG formula disk}
{\rm det}_\zeta(\Lambda_0)={\rm det}_\zeta(\partial_\gamma)=|\Gamma|.
\end{equation}

Introduce the {\it Hilbert transform} of $M$ by
\begin{equation*}
%\label{Hilbert}
H:=\partial_\gamma^{-1}\Lambda.
\end{equation*}
It is easy to see that the operator $f+\mathbb{C}\mapsto Hf+\mathbb{C}$ acting on the quotient space $L_2(\Gamma,dl_g;\mathbb{C})/\mathbb{C}$ depends only on the complex structure on $M$ (if $H'$ is the Hilbert transfrom of $(M,h:=e^\phi g)$, then $H'f=Hf+c(f)$, where $c(f)$ is the constant providing the orthogonality of $H'f$ to the constants in $L_2(\Gamma,dl_h;\mathbb{C})$). Conversely, the Hilbert transforms $H$ determines the surface $M$ up to biholomorphism \cite{B}. 

Decomposition (\ref{DN PDO decomp}) implies
\begin{equation}
\label{Hilbert PDO decomp}
H=H_0+\tilde{H}, \qquad H_0=\partial_\gamma^{-1}|\partial_\gamma|,
\end{equation}
where $\tilde{H}$ is a smoothing $\Psi$DO and $H_0$ is the Hilbert transform of $\mathbb{D}$. Note that $H_0^2=-I$. The spectrum of $H$ is described in \cite{HilbKor}. The essential spectrum of $H$ consists of two eigenvalues $-i,+i$; the corresponding eigenspaces ${\rm Ker}(H+i)$ (${\rm Ker}(H-i)$) are spanned by the boundary traces of holomorphic (anti-holomorphic) functions with square-integrable differentials on $M$. The discrete spectrum of $H$ consists of the eigenvalues (having the same algebraic and geometric finite multiplicities)
$$\lambda_{\pm 1}(H)=\pm i\mu_1,\dots,\lambda_{\pm\mathfrak{g}}(H)=\pm i\mu_{\mathfrak{g}}$$ 
(here the eigenvalues are repeated according to their multiplicities). The numbers $\mu_k\in (0,1)$ are related to the periods of Abelian differentials on the Schottky double $2M$ of $M$ (equipped with the antiholomorphic involution $\tau$, $2M/\tau\equiv M$. 
\begin{lemma}[see Corollary 4, \cite{znsbKor}] \label{Kolemma} There is the basis $\{\omega_{k},\omega_{-k}\}_{k=1}^{\mathfrak{g}}$ in $H^0(2M;K)$ obeying
$$\int_{\tau\circ l}\omega_{\pm k}=\frac{\pm \mu_k+1}{\pm\mu_k-1}\int_{l}\omega_{\pm k}$$
for any loop $l$ in $M\subset 2M$.
\end{lemma}
(For the convenience of the reader, the proof of this statement is presented in the appendix. Also, we deduce equation (\ref{finding mus}) expressing $\lambda_{\pm k}$ in terms of the b-period matrix in symmetric homology basis on $2M$.)

Due to the presence of the essential spectrum $\{-i,+i\}$, neither zeta-regularized nor Fredholm determinant of $H$ can be defined.We simply define  the determinant of $H$ via
\begin{equation}
\label{toy det}
{\rm DET}(H):=\lambda_1\cdot\dots\cdot\lambda_{\mathfrak{g}}\cdot\lambda_{-1}\cdot\dots\cdot\lambda_{-\mathfrak{g}}={\rm det}(H_{\rm disc}).
\end{equation}
i. e. by restricting $H$ to its discrete spectrum subspace $\mathcal{H}_{\rm disc}:=\bigoplus\limits_{\pm k}{\rm Ker}(H-\lambda_k)$ of dimension $2\mathfrak{g}$; such a restriction is denoted by $H_{\rm disc}$ (thus, roughly speaking, we formally put $(-i)(+i)(-i)(+i)\dots=1$ for the eigenvalues of $H$ of infinite multiplicity, by this definition, we have ${\rm DET}(H_0)=1$).

\smallskip

In this note, we prove the following relation 
\begin{equation}
\label{dumb formula}
{\rm det}_\zeta(\Lambda)={\rm det}_\zeta(\partial_\gamma)\cdot{\rm DET}(H).
\end{equation}
As a corollary of (\ref{dumb formula}), (\ref{GG formula disk}), (\ref{toy det}), and (\ref{GG formula}), one arrives to
\begin{align*}
(\mu_1\cdot\dots\cdot\mu_{\mathfrak{g}})^2={\rm DET}(H)={\rm det}_\zeta(\Lambda)/{\rm det}_\zeta(\partial_\gamma)=\\
={\rm det}_\zeta(\Lambda)/|\Gamma|=\Big((2\pi s)^{-\mathfrak{g}}\mathscr{R}(s)\Big)\Big|_{s=0}/(1-\mathfrak{g}).
\end{align*}
Thus, we have obtained the formula
\begin{equation}
\label{KoKoFormula}
\Big((2\pi s)^{-\mathfrak{g}}\mathscr{R}(s)\Big)\Big|_{s=0}=(1-\mathfrak{g})(\mu_1\cdot\dots\cdot\mu_{\mathfrak{g}})^2
\end{equation}
that relates the length spectrum on $(M\backslash\Gamma,h_\infty)$ with the periods of certain basic Abelian differentials on the double $(2M,\tau)$. 

\section*{Proof of formula (\ref{dumb formula})}
The proof of (\ref{dumb formula}) is based on the Kontsevich-Vishik definition of the determinant of $\Psi$DO and the  variational approach of \cite{Ww,F}.

{\bf i)} First, recall the basic facts and definitions. Let $A$ be an invertible $\Psi$DO of order $l\ge 0$. We assume that $A$ is sectorial, i.e. its spectrum is contained in some closed sector $\mathbb{K}:=\{\lambda\in\mathbb{C} \ | \ {\rm arg}(\lambda)\in[\varphi_1,\varphi_2]\}\ne\mathbb{C}$ and $\|A-\lambda I\|_{B(\mathcal{H})}=O(1/{\rm dist}(\Lambda,\mathbb{K}))$. Then one can define ${\rm log}(A)$ by using the basic holomorphic calculus
$${\rm log}(A):=\frac{1}{2\pi i}\int\limits_{\mathfrak{C}}A(\lambda I-A)^{-1}\frac{{\rm log}(\lambda)}{\lambda}d\lambda,$$
where $\mathfrak{C}$ is the (counterclockwise) curve enclosed $Sp(A)$ in such a way that $\|A-\lambda I\|_{B(\mathcal{H})}=O(1/|\lambda|)$ (say, $\mathfrak{C}$ is bounded or coincides with the rays ${\rm arg}(\lambda)=\varphi_1-\epsilon$, ${\rm arg}(\lambda)=\varphi_2+\epsilon$ ($\epsilon>0$) near infinity) for $\lambda\in\mathfrak{C}$, the cut of the logarithm is a ray intersecting $\mathbb{K}$ only at zero. 

Introduce some regularizer $Q$ which is a positive 1st order elliptic $\Psi$DO (as such $Q$, one can take $\Lambda$ or $\Lambda_0$). The {\it Kontsevich–Vishik determinant} of $A$ is defined as
\begin{equation*}
%\label{KZ det}
{\rm det}_Q(A):={\rm exp}\Bigg(\underset{s=0}{\rm f.p.}\,\,\Big[a.c.\,\,{\rm Tr}\Big(Q^{-s}{\rm log}(A)\Big)\Big]\Bigg),
\end{equation*}
where $a.c.$ is the analytic continuation from the half-plane $\Re s>s_0\gg 1$ in which $Q^{-s}{\rm log}(A)$ is of trace class, and ${\rm f.p.}$ means the Hadamard finite part. 

Next, let $[0,1]\ni t\mapsto A(t)$ be a differentiable family of operators obeying the above conditions, $\dot{A}(t)=\partial_t A(t)$ is of trace class for all $t$, and $t\mapsto \dot{A}(t)$ is continuous in the trace-norm $\|\cdot\|_1$. We have
\begin{equation}
\label{variation of logdet}
\partial_t{\rm log}\,{\rm det}_Q(A)=\frac{1}{2\pi i}\underset{s=0}{\rm f.p.}\,\,\Big[a.c.\,\,\Big(\int\limits_{\mathbb{C}}{\rm Tr}\Big(Q^{-s}\partial_t(A(\lambda I-A)^{-1})\Big)\frac{{\rm log}(\lambda)}{\lambda}d\lambda\Big]
\end{equation}
(if the right-hand is well-defined for $\Re s>-\epsilon$ ($\epsilon>0$), then one ensures that one can interchange the differentiation in $t$ and the analytic continuation in $s$ in the above formula by integrating it in $t$). We have
$$\partial_t(A(\lambda I-A)^{-1})=\dot{A}(z-A)^{-1}+(z-A)^{-1}A\dot{A}(z-A)^{-1},$$
Note that $(z-A)^{-1}A=-I+z(z-A)^{-1}$ is bounded and thus $Q^{-s}(z-A)^{-1}A\dot{A}$ is of trace class. Hence,
$${\rm Tr}\Big(Q^{-s}\partial_t(A(\lambda I-A)^{-1})\Big)={\rm Tr}\Big((z-A)^{-1}Q^{-s}\dot{A}+(z-A)^{-1}Q^{-s}(z-A)^{-1}A\dot{A}\Big)$$
and the well-known estimate $|{\rm Tr}(B'B)|\le\|B'\|_{1}\|B\|_{B(H)}$ yields
\begin{align*}
\lim_{s\to 0}{\rm Tr}\Big(Q^{-s}\partial_t(A(\lambda I-A)^{-1})\Big)={\rm Tr}\Big(((z-A)^{-1}+(z-A)^{-2}A)\dot{A}\Big)=\\
={\rm Tr}\Big(\Big((z-A)^{-1}-\partial_z\big(z(z-A)^{-1}\big)\Big)\dot{A}\Big).
\end{align*}
and the integration by parts in (\ref{variation of logdet}) leads to
\begin{align}
\label{variation formula}
\begin{split}
\partial_t{\rm log}\,{\rm det}_Q&(A)=\\
={\rm Tr}\Bigg(\frac{1}{2\pi i}&\int\limits_{\mathfrak{C}}\Bigg(\frac{{\rm log}(z)}{z}+z\Big(\frac{1}{z^{2}}-\frac{{\rm log}(z)}{z^2}\Big)\Bigg)(z-A)^{-1}dz\dot{A}\Bigg)={\rm Tr}(A^{-1}\dot{A})
\end{split}
\end{align}
(boundary terms are of order $z\|(z-A)^{-1}\|_{B(\mathcal{H})}{\rm log}(z)/z=O(|{\rm log}(z)/z|)$ and this vanish at infinity).

For the two families $A$ and $B$ obeying the above conditions, the $\partial_t(AB)=\dot{A}B+A\dot{B}$ is trace class and is continuous in $t$ in the trace norm. Then, as immediate corollary of (\ref{variation formula}), one obtains
\begin{align*}
\partial_t{\rm log}\,{\rm det}_Q&(AB)={\rm Tr}\big((AB)^{-1}\partial_t(AB)\big)={\rm Tr}\big(B^{-1}A^{-1}\dot{A}B+B^{-1}\dot{B}\big)=\\
=&{\rm Tr}\big(B^{-1}A^{-1}\dot{A}B\big)+\partial_t{\rm log}\,{\rm det}_Q(B)= \ \text{($\dot{A}B$ is of trace class)}\\
=&{\rm Tr}\big(\dot{A}BB^{-1}A^{-1}\big)+\partial_t{\rm log}\,{\rm det}_Q(B)=\partial_t({\rm log}\,{\rm det}_Q(A)+{\rm log}\,{\rm det}_Q(B)),
\end{align*}
i.e. the following multiplicativity formula is valid
\begin{align}
\label{multiplicativity deformation}
\frac{{\rm det}_Q(A(1)B(1))}{{\rm det}_Q(A(1)){\rm det}_Q(B(1))}=\frac{{\rm det}_Q(A(0)B(0))}{{\rm det}_Q(A(0)){\rm det}_Q(B(0))}
\end{align}
regardless of the paths connecting $A(0)$ ($B(0)$) and $A(1)$ ($B(1)$) (as long as they satisfy the above conditions).

{\bf ii)} Consider the families 
\begin{align*}
A(t)=\partial_\gamma, \quad B(t)=\theta(t)H+(1-\theta(t))H_0=H_0+\theta(t)\tilde{H}, \\ 
A(t)B(t)=\theta(t)\Lambda+(1-\theta(t))\Lambda_0=\Lambda_0+\theta(t)\tilde{\Lambda},
\end{align*}
where $\theta$ is the smooth curve in $\mathbb{C}$ connecting $\theta(0)=0$ and $\theta(1)=1$. One can chose $\theta$ in such a way that $B(t)$, $A(t)B(t)$ are invertible for any $t$. Indeed, $B(\theta)=H_0+\theta\tilde{H}=H_0^{-1}(\theta H_0\tilde{H}-I)$ is non-invertible if and only if $\theta^{-1}$ belongs to the spectrum of $H_0\tilde{H}$; since $H_0\tilde{H}$ is a compact operator (moreover, is a smoothing $\Psi$DO), its spectrum has only one accumulation point $\theta^{-1}=0$. In addition, the derivatives
$$\dot{A}=0, \quad \dot{B}=\dot{\theta}\tilde{H}$$
are of trace class. So, the families $A,B$ obeys the condition of {\bf i)} and thus formula (\ref{multiplicativity deformation}) is valid for them and implies
\begin{align}
\label{multiplicativity DN Hilbert diff}
\frac{{\rm det}_Q(\Lambda)}{{\rm det}_Q(\partial_\gamma){\rm det}_Q(H)}=\frac{{\rm det}_Q(\Lambda_0)}{{\rm det}_Q(\partial_\gamma){\rm det}_Q(H_0)}.
\end{align}

{\bf iii)} Introduce the family 
$$J(t):=-B(t)^2, \qquad B(t)=H_0+\theta(t)\tilde{H}.$$
It derivative $\dot{J}(t)=-\dot{\theta}(t)(B(t)\tilde{H}+\tilde{H}B(t))$ is a smoothing operator. Thus, repeating of the above reasoning yields
\begin{align*}
\partial_t{\rm log}\,{\rm det}_Q(J)={\rm Tr}(J^{-1}\dot{J})=\dot{\theta}(t){\rm Tr}_Q\Big(B(t)^{-1}\tilde{H}+B(t)^{-2}\cdot\tilde{H}B(t)\Big)=\\
=2\dot{\theta}(t){\rm Tr}_Q\Big(B(t)^{-1}\tilde{H}\Big)=2\partial_t{\rm log}\,{\rm det}_Q(B).
\end{align*}
Integrating and taking into account that $B(0)=H_0$, $B(1)=H$, $-H_0^2=I$, and ${\rm det}_Q(I)=0$, one obtains
\begin{align*}
{\rm det}_Q(-H^2)=&\frac{{\rm det}_Q(-H^2)}{{\rm det}_Q(I)}=\frac{{\rm det}_Q(-H^2)}{{\rm det}_Q(-H_0^2)}=\\
=&\frac{{\rm det}_Q(J(1))}{{\rm det}_Q(J(0))}=\Big(\frac{{\rm det}_Q(B(1))}{{\rm det}_Q(B(0))}\Big)^2=\Big(\frac{{\rm det}_Q(H)}{{\rm det}_Q(H_0)}\Big)^2.
\end{align*}
At the same time, decomposition (\ref{Hilbert PDO decomp}) implies that $K:=I+H^2$ is a smoothing $\Psi$DO, whence
\begin{align*}
{\rm log}{\rm det}_Q(-H^2)={\rm Tr}\big(Q^{-s}{\rm log}(I-K)\big)\big|_{s=0}=\sum_{k=1}^{\infty}\frac{{\rm Tr}\big(K^k\big)}{k}.
\end{align*}
Note that ${\rm Ker}(K)=\mathcal{H}\ominus\mathcal{H}_{\rm disc}$ and $K_{\rm disc}:=K|_{\mathcal{H}_{\rm disc}}=I+H_{\rm disc}^2$. Then the above formula implies
\begin{align*}
{\rm log}{\rm det}_Q(-H^2)=\sum_{k=1}^{\infty}\frac{{\rm Tr}\big(K_{\rm disc}^k\big)}{k}={\rm Tr}\big({\rm log}(-H_{\rm disc}^2)\big)=\\
=(-1)^{{\rm dim}(\mathcal{H}_{\rm disc})}{\rm log}({\rm det}(H_{\rm disc})^2)={\rm log}({\rm DET}(H)^2).
\end{align*}
Thus, we finally obtain
\begin{equation}
\label{Q det of Hilbert}
\frac{{\rm det}_Q(H)}{{\rm det}_Q(H_0)}={\rm DET}(H).
\end{equation}
Now combining (\ref{multiplicativity DN Hilbert diff}) with (\ref{Q det of Hilbert}) yields
\begin{align}
\label{Q-multiplicativity}
{\rm det}_Q(\Lambda)={\rm det}_Q(\Lambda_0)\cdot{\rm DET}(H).
\end{align}

{\bf iv)} Let us show that ${\rm det}_{Q}(\Lambda)$ is independent of the choice of the regularizer $Q=\Lambda,\Lambda_0$. To this end, we let us consider the family of regularizers 
$$t\mapsto Q(t)=\Lambda_0+\alpha(t)\tilde{\Lambda}=\Lambda_0+\alpha(t)\partial_\gamma\tilde{H},$$ 
where the second term is again a smoothing $\Psi$DO, and $\alpha$ is again a smooth curve connecting $0$ and $1$ in such a way that $Q(t)$ is invertible for all $t$. Then 
$$\dot{Q}=\dot{\alpha}\partial_\gamma\tilde{H}$$
is a smoothing $\Psi$DO. We have
$$\partial_t{\rm Tr}(Q^{-s}{\rm log}\Lambda)=-s{\rm Tr}(Q^{-s-1}\dot{Q}{\rm log}\Lambda);$$
here the trace in the right-hand side is well defined for all $s$, so one can put $s=0$ and obtain $\partial_t{\rm Tr}(Q_t^{-s}{\rm log}\Lambda)|_{s=0}=0$. Thus, we have 
\begin{align}
\label{regularizer independence}
{\rm log}\,{\rm det}_{\Lambda_0}(\Lambda)={\rm Tr}(Q(0)^{-s}{\rm log}\Lambda)={\rm Tr}(Q(1)^{-s}{\rm log}\Lambda)={\rm log}\,{\rm det}_{\Lambda}(\Lambda).
\end{align}

{\bf v)} Choosing $Q=\Lambda_0$ or $Q=\Lambda$ as a regularizer, one arrives at
\begin{align*}
{\rm log}\,{\rm det}_{\Lambda}(\Lambda)=a.c.{\rm Tr}(\Lambda^{-s}{\rm log}(\Lambda))\big|_{s=0}=a.c.\big[-\partial_s{\rm Tr}(\Lambda^{-s})\big]\big|_{s=0}=\\
=-\zeta'_{\Lambda}(0)={\rm log}\,{\rm det}_\zeta(\Lambda),
\end{align*}
whence
\begin{equation}
\label{KZ and zeta det connection}
{\rm det}_{\Lambda}(\Lambda)={\rm det}_\zeta(\Lambda), \qquad {\rm det}_{\Lambda_0}(\Lambda_0)={\rm det}_\zeta(\Lambda_0)={\rm det}_\zeta(\partial_\gamma)
\end{equation}
(the last equality follows from (\ref{GG formula disk})). Now substitution of (\ref{regularizer independence}) and (\ref{KZ and zeta det connection}) into (\ref{Q-multiplicativity}) yields
\begin{align*}
{\rm det}_\zeta(\Lambda)={\rm det}_{\Lambda}(\Lambda)={\rm det}_{\Lambda_0}(\Lambda)={\rm det}_{\Lambda_0}(\Lambda_0)\cdot{\rm DET}(H)={\rm det}_\zeta(\partial_\gamma)\cdot{\rm DET}(H).
\end{align*}
By this formula (\ref{dumb formula}) is proved. \qed

\section{Appendix: proof of Lemma \ref{Kolemma}}
Suppose that $Hf=-\lambda f\ne 0$, i.e. $\Lambda f=-\lambda\partial_\gamma f$ on $\Gamma$ and denote by $u^f$ the harmonic extension of $f$ into $M$. Let $\star$ denotes the Hodge operator on $2M\supset M$ and $\gamma$ is the unit tangent vector obtained by the counterclockwise rotation on the right angle of $\nu$. The space of harmonic 1-forms $\mathscr{H}^1=\{a\in L_2(M;T^*M) \ | \ da=d\star a = 0 \text{ in } M\}$ admits the orthogonal decomposition $\mathscr{H}^1=d\mathscr{H}^0\oplus\mathscr{N}$, where $\mathscr{H}^0=\mathscr{H}^1\cap dH^1(M;\mathbb{C})$ is the sub-space of the exact forms and $\mathscr{N}:=\{a\in \mathscr{H}^1 \ | \ a(\nu)=0\}$ is the sub-space of harmonic fields tangent to $\Gamma$. Thus, we have
$$\star du^f=du^h+p, \qquad \star du^h=du^r+q \qquad (p,q\in\mathscr{N}).$$
In particular, $q+\star p=-du^{r+f}$. Restricting the above equations to $\Gamma$, one arrives at
$$\Lambda r=\star(du^h-q)(\nu)=-\partial_\gamma h, \qquad \Lambda h=-\partial_\gamma f=\lambda^{-1}\Lambda f.$$
Hence $h=\lambda^{-1}f+{\rm const}$, $\Lambda r=-\lambda^{-1}\partial_\gamma f=\lambda^{-2}\Lambda f$, and $r=\lambda^{-2}f+{\rm const}$. So, we have 
$$q+\star p=-(\lambda^{-2}+1)du^{f},$$ 
in particular, 
$$p(\gamma)=-[q+\star p](\nu)=(\lambda^{-2}+1)\Lambda f=-(\lambda^{-2}+1)\lambda\partial_\gamma f=\lambda[q+\star p](\gamma)=\lambda q(\gamma),$$
whence $p=\lambda q$ due to the uniqueness of solutions to the Cauchy problem for elliptic equations (Tataru-Holmgren-Kovalevskaya theorem). 

So, we have proved that the form $\tilde{\omega}:=q+\lambda\star q$ is exact. Since $q$ is tangent to $\Gamma$, it admits the harmonic extension (still denoted by $q$) onto the double $2M\supset M$ by the rule $\tau^*q=q$. Introduce the Abelian differrential 
$$\omega:=(i-\star)q;$$
note that its trace on $\Gamma$ obeys
\begin{equation}
\label{linear independence}
\omega(\gamma)=\omega(\nu)=iq(\gamma)=-(\lambda^{-2}+1)\partial_\gamma f.
\end{equation}
Let $ l$ be a closed curve in $M$. Since the involution on $2M$ is anti-holomorphic, $\tau^*\star=-\star\tau^*$, one has
\begin{align*}
\int\limits_{\tau\circ l}\omega=&\int\limits_{ l}\tau^*\omega=\int\limits_{ l}\tau^*(i-\star)q=\int\limits_{ l}(i+\star)\tau^*q=\\=&\int\limits_{ l}(i+\star)q=\int\limits_{ l}(1-\lambda)\star q+\int\limits_{ l}i\tilde{\omega}=(1-i\lambda)\int\limits_{ l}\star q+0.
\end{align*}
At the same time,
\begin{equation}
\label{eigenad}
\int\limits_{ l}\omega=\int\limits_{ l}(i\tilde{\omega}-(1+i\lambda)\star q)=-(1+i\lambda)\int\limits_{ l}\star q+0.
\end{equation}
Comparing the last two equalities, one obtains
$$\int\limits_{\tau\circ l}\omega=\frac{\lambda+i}{\lambda-i}\int\limits_{ l}\omega.$$
By this, Lemma \ref{Kolemma} is proved.\qed

\

Let $\{a_i,b_i\}_{i=1}^{\mathfrak{2g}}$ be a homology basis on the double $(2M,\tau)$ obeying symmetry conditions $a_{\mathfrak{g}+i}=\tau\circ a_{i}$, $b_{\mathfrak{g}+i}=-\tau\circ b_{i}$ and such that $\{a_i,b_i\}_{i=1}^{\mathfrak{g}}$ can be represented by the loops contained in $M$; then the dual basis $\{\theta_i\}_{i=1}^{2\mathfrak{g}}\subset H^0(2M;K)$, $\int_{a_i}\theta_j=\delta_{ij}$ satisfies $\theta_{\mathfrak{g}+i}=\overline{\theta_{i}\circ\tau}$. Hence, the corresponding $b$-period matrix $\mathbb{B}$ is of the form
\begin{align*}
\mathbb{B}=i\left(\begin{array}{cc}
\mathfrak{G} & \mathfrak{F}\\
\overline{\mathfrak{F}} & \overline{\mathfrak{G}}
\end{array}\right) \qquad (\mathfrak{F}^*=\mathfrak{F}, \ \mathfrak{G}^T=\mathfrak{G}).
\end{align*}
Condition (\ref{eigenad}) for $\omega=\sum_{k=1}^{2\mathfrak{g}}c_k\theta_k$ is equivalent to equations
\begin{align*}
\frac{c_{\mathfrak{g}+j}}{c_j}=\frac{\lambda+i}{\lambda-i}, \qquad \frac{\sum_{k=1}^{2\mathfrak{g}}\mathbb{B}_{\mathfrak{g}+j,k}c_k}{\sum_{k=1}^{2\mathfrak{g}}\mathbb{B}_{j,k}c_k}=-\frac{\lambda+i}{\lambda-i} \qquad (j=1,\dots,\mathfrak{g}).
\end{align*}
which are compatible if and only if
\begin{equation}
\label{finding mus}
{\rm det}\Bigg(\frac{\lambda+i}{\lambda-i}\cdot\mathfrak{F}+2\Re\mathfrak{G}+\frac{\lambda-i}{\lambda+i}\cdot\overline{\mathfrak{F}}\Bigg)=0.
\end{equation}
Thus, the solutions to (\ref{finding mus}) are exactly the eigenvalues of $H_{\rm disc}$.

\begin{thebibliography}{99}
\bibitem{EV}
\newblock{J. Edward, S. Wu} 
\newblock{\em  Determinant of the Neumann operator on smooth Jordan curves.}
\newblock{Proc. AMS, 111(2), 357--363 (1991).}

\bibitem{Guillarmou} 
\newblock{Colin Guillarmou, Laurent Guillop\'e}
\newblock{\em  The determinant of the Dirichlet-to-Neumann map for surfaces with boundary.}
\newblock{(2007), \url{https://arxiv.org/abs/math/0701727}}

\bibitem{B}
\newblock{M.I. Belishev}
\newblock{\em  The Calderon problem for two-dimensional manifolds by the BC-method.} 
\newblock{SIAM Journal of Mathematical Analysis 35(1), 172--182 (2003). \url{https://doi.org/10.1137/S0036141002413919}}

\bibitem{HilbKor} 
\newblock{Dmitrii Korikov}
\newblock{\em  Determination of period matrix of double of surface with boundary via its DN map.}
\newblock{Canadian Journal of Mathematics (2025), \url{https://doi.org/10.4153/s0008414x25000264}}

\bibitem{KV} 
\newblock{Maxim Kontsevich, Simeon Vishik}
\newblock{\em  Determinants of elliptic pseudo-differential operators.}
\newblock{(1994), \url{https://arxiv.org/pdf/hep-th/9404046}}

\bibitem{Ww} 
\newblock{Richard A. Wentworth}
\newblock{\em  Gluing formulas for determinants of Dolbeault laplacians on Riemann surfaces.}
\newblock{(2012), \url{https://arxiv.org/abs/1008.2914}}

\bibitem{F} 
\newblock{Leonid Friedlander}
\newblock{\em  On Multiplicative Properties of Determinants.}
\newblock{(2018), \url{https://doi.org/10.48550/arXiv.1801.10606}}

\bibitem{Grubb} 

\newblock{G.Grubb}
\newblock{\em  Functional calculus of pseudo-differential boundary value problems.}
\newblock{Birkhaeuser, Boston, 1986}

\bibitem{LeeU}
\newblock{J.M. Lee, G. Uhlmann} 
\newblock{\em  Determining anisotropic real‐analytic conductivities by boundary measurements.}
\newblock{Comm. Pure Appl. Math., 42, 1097--1112 (1989).}

\bibitem{znsbKor} 
\newblock{Dmitrii Korikov}
\newblock{\em  On non-homeomorphic surfaces with close DN maps.}
\newblock{(2026), \url{https://arxiv.org/abs/2602.13236}}

\end{thebibliography}


\end{document}

